% Propietario conservado: /Users/ruben/Documents/New project/output/REVISION_ARGUMENTAL_APERTURAS_CIERRES_20260915/VII/delivery/MANUSCRIPT_VII_ES_EN_COMPLETE/source_es/manuscrito/28_portador_tensorial.tex
% FORMALIZACION_REUNIDA: representación pentacomponente de II,
% con antecedente Paley y demostración directa de los momentos del marco.
% El retorno transporta la carta; no se identifica un ciclo de orden doce con A5.
\section{Oriented incidence and gravitational tensor representation}
\label{vii:vii:sec:pentacomponente}
The exceptional publication of the state preserves the incidence accompanying
its arithmetic reading. The ternary matrix \(A_W\) of
\eqref{vii:exc:aw} provides a precise material antecedent here. Its
balanced lift \(C=\operatorname{bal}(A_W)\), with representatives
\(0,1,-1\), satisfies
\[
 C=C^{\mathsf T},\qquad C^2=5I_6,\qquad \operatorname{Tr}C=0.
\]
The six positions have the order \((\infty,0,1,2,3,4)\).
This incidence first produces a three-dimensional space; its quadratic
reading then produces the five-component space on which
the gravitational coupling acts.

\subsection{Signed covering and calendar transport}
We define
\[
 E_C=\frac12(I_6+C/\sqrt5),\quad
 \mathcal H_C=\operatorname{im}E_C,\quad
 v_{u,\sigma}=\sigma\sqrt2E_Ce_u,\quad \sigma\in\{1,-1\}.
\]
\begin{proposition}[Twelve oriented positions]
\label{vii:vii:prop:doce-orientadas}
The vectors \(v_{u,\sigma}\) are twelve distinct unit vectors
in a three-dimensional space. Their incidence is determined by
\(\langle v_{u,\sigma},v_{v,\tau}\rangle=1/\sqrt5\).
Orientation reversal acts by antipodality.
\end{proposition}
\begin{proof}
The quadratic identity for \(C\) implies \(E_C^2=E_C=E_C^*\);
its trace is three. Since the diagonal of \(C\) is zero,
\[
 \langle v_{u,\sigma},v_{v,\tau}\rangle
 =\begin{cases}\sigma\tau,&u=v,\\
 \sigma\tau C_{uv}/\sqrt5,&u\ne v.
 \end{cases}
\]
This proves normalization, distinction of the twelve vectors and incidence.
Changing \(\sigma\) to \(-\sigma\) changes the sign of the vector.
\end{proof}

An explicit Euclidean expression is obtained by setting
\(\nu=\sqrt{1+\varphi^2}\) and
\[
 B=\frac1\nu
 \begin{pmatrix}
 0&-1&-\varphi&0&\varphi&1\\
 -1&-\varphi&0&1&0&-\varphi\\
 -\varphi&0&-1&-\varphi&-1&0
 \end{pmatrix}.
\]
The relation \(\varphi^2=\varphi+1\), for the coordinate already generated,
gives \(B^*B=2E_C\).
Consequently \(B/\sqrt2|_{\mathcal H_C}\) is an isometry and takes
\(v_{u,\sigma}\) to \(\sigma Be_u\). The image is
\[
 \mathcal V=\frac1\nu\bigl(
 \{0\}\times\{\pm1\}\times\{\pm\varphi\}\ \cup\
 \{\pm1\}\times\{\pm\varphi\}\times\{0\}\ \cup\
 \{\pm\varphi\}\times\{0\}\times\{\pm1\}\bigr).
\]
The internal coordinate, the incidence and its
twelve geometric representatives are thus related.

The calendar coordinate is written \((b,r)\in
\mathbb Z/12\mathbb Z\times\{0,\ldots,8\}\). A route of length \(n\)
induces
\[
 k(r,n)=\left\lfloor\frac{r+n}{9}\right\rfloor,\qquad
 (b,r)\longmapsto
 \bigl(b+k(r,n)\bmod12,\ (r+n)\bmod9\bigr).
\]
The explicit alignment of its positions \(p=b+1\) with the signed chart is
\[
\begin{array}{c|rrrrrrrrrrrr}
p&1&2&3&4&5&6&7&8&9&10&11&12\\ \hline
u&\infty&0&0&\infty&1&3&3&1&2&4&4&2\\
\sigma&+&-&+&-&-&+&-&+&-&+&-&+
\end{array}
\]
This table fixes a trivialization of the carrier. Each change of position
jointly transports the vectors and their incidence matrix.

\begin{proposition}[Covariance of the frame under return]
If \(S e_p=e_{p+1\bmod12}\), the matrix representation of an
operator \(P\) of the carrier is transported as \(P_k=S^kPS^{-k}\).
This transport preserves inner products, incidences and route
composition. After twelve returns the chart returns, preserving the memory
accumulated in the source state.
\end{proposition}
\begin{proof}
\(S\) is orthogonal and \(S^{12}=I\); the metric properties
and periodicity of the chart follow. For composition, with
\(r'=(r+n)\bmod9\), Euclidean division gives
\(k(r,n+m)=k(r,n)+k(r',m)\). The exponents of \(S\) are added
exactly with the carry. The operation acts on the calendar
coordinate; the route's memory updates remain
present in its enriched domain.
\end{proof}

\subsection{Quadratic reading and five-component space}
Let \(\mathbb R^{\mathcal V}\) be the space of the twelve coordinates,
\(\mathsf A e_v=e_{-v}\) its antipodal involution and
\(\mathsf J=\mathbf1\mathbf1^{\mathsf T}\). We define
\[
 P_5=\frac12(I+\mathsf A)-\frac1{12}\mathsf J,\qquad
 Q_v=vv^{\mathsf T}-\frac13I_3,\qquad
 \Gamma_0x=\sum_{v\in\mathcal V}x_vQ_v.
\]
\begin{theorem}[Five-component projector and tensor isometry]
\label{vii:vii:thm:portador-pentacomponente}
\(P_5\) is the rank-five orthogonal projector onto the space of
coordinates even under antipodality and of zero sum. Moreover,
\[
 \Gamma_0^*\Gamma_0=\frac85P_5,\qquad
 \Gamma_0\Gamma_0^*=\frac85I_{\operatorname{Sym}^2_0(\mathbb R^3)}.
\]
Therefore
\[
 \Gamma_5=\sqrt{\frac58}\,\Gamma_0|_{\operatorname{im}P_5}
\]
is an isometry onto the space of traceless symmetric tensors.
\end{theorem}
\begin{proof}
The even space has one coordinate for each antipodal pair, so
its dimension is six. The constant vector belongs to it; subtracting its
rank-one projector gives \(P_5\), of rank five.

Since
\(\langle Q_v,Q_w\rangle_F=(v^{\mathsf T}w)^2-1/3\), the entries
of \(\Gamma_0^*\Gamma_0\) equal \(2/3\) for \(w=\pm v\) and
\(-2/15\) in the other ten positions of each row. These are
exactly the entries of \((8/5)P_5\).

The identity can also be checked in tensor space without
resorting to a classification of representations. The list of vertices gives
\[
 \sum_vv_i^4=\frac{12}{5},\qquad
 \sum_vv_i^2v_j^2=\frac45\quad(i\ne j),
\]
and moments with any odd power vanish by sign
changes. Indeed,
\((1+\varphi^2)^2=5\varphi^2\) and \(1+\varphi^4=3\varphi^2\).
For a symmetric tensor \(Q\), these sums give
\[
 \sum_v(v^{\mathsf T}Qv)\,vv^{\mathsf T}
 =\frac85Q+\frac45(\operatorname{tr}Q)I.
\]
With \(\operatorname{tr}Q=0\), and using
\(\sum_vv^{\mathsf T}Qv=4\operatorname{tr}Q=0\), we obtain
\(\Gamma_0\Gamma_0^*Q=(8/5)Q\). The two identities prove
the normalized isometry and its surjectivity.
\end{proof}

The gravitational section in \(\ref{vii:v:ant:sec:gravity}\) determines the
operator \(G_aP_5\). Its spectrum has \(G_a\) with multiplicity
five and zero with multiplicity seven. Orientation is preserved in
the carrier and its transport; \(G_a\) fixes the common coupling.

\subsection{Transverse projection and helicities}
For \(n\in S^2\), set \(\Pi=I-nn^{\mathsf T}\) and
\[
 \Lambda_nQ=\Pi Q\Pi-\frac12\operatorname{tr}(\Pi Q\Pi)\Pi.
\]
\begin{proposition}[Rank-two transverse reduction]
\(\Lambda_n\) is the orthogonal projector onto
\(\{Q\in\operatorname{Sym}^2_0(\mathbb R^3):Qn=0\}\), of dimension two.
\end{proposition}
\begin{proof}
\(\Pi n=0\) gives transversality and \(\operatorname{tr}\Pi=2\) gives
zero trace. A transverse traceless tensor remains fixed.
The formula is self-adjoint for the Frobenius inner product.
In a frame \((e_1,e_2,n)\), the image has orthonormal basis
\[
 E_+=\frac{e_1e_1^{\mathsf T}-e_2e_2^{\mathsf T}}{\sqrt2},\qquad
 E_\times=\frac{e_1e_2^{\mathsf T}+e_2e_1^{\mathsf T}}{\sqrt2}.
\]
\end{proof}
The other three vectors of an orthonormal basis of the carrier are
\[
 E_1=\frac{e_1n^{\mathsf T}+ne_1^{\mathsf T}}{\sqrt2},\quad
 E_2=\frac{e_2n^{\mathsf T}+ne_2^{\mathsf T}}{\sqrt2},\quad
 E_0=\frac{2nn^{\mathsf T}-e_1e_1^{\mathsf T}-e_2e_2^{\mathsf T}}{\sqrt6}.
\]
When the transverse frame is rotated by an angle \(\theta\), the pairs
\((E_+,E_\times)\) and \((E_1,E_2)\) rotate by \(2\theta\) and
\(\theta\), respectively, while \(E_0\) remains fixed. The
complexification thus has weights \(2,-2,1,-1,0\).
\(\Lambda_n\) preserves the first two and annihilates the rest.
The chain is established
\[
 \mathbb R^{12}\xrightarrow{P_5}\operatorname{im}P_5
 \xrightarrow{\Gamma_5}\operatorname{Sym}^2_0(\mathbb R^3)
 \xrightarrow{\Lambda_n}\mathcal H_{\rm TT}(n),
 \qquad 12\longrightarrow5\longrightarrow5\longrightarrow2.
\]
The five-component representation and its reduction are algebraic
results preceding the choice of field equations.
In the massless gravitational realization, the gauge constraints
select the transverse image. This distinction preserves both the
five structural coordinates and the two radiative helicities.

